% !TeX program = xelatex
% UTF-8; standalone XeLaTeX/XePersian document, or input in an existing XePersian document.
% Complete transcription of Google Drive: Notebook/brain tumors.
% Original photographs are saved unchanged in neurology/pbrain tumors/images:
% pic01.jpg = 20261012_210533558.jpg
% pic02.jpg = 20261012_210540633.jpg
% pic03.jpg = 20261012_210549787.jpg
% pic04.jpg = 20261012_210604562.jpg
% pic05.jpg = 20261012_210611285.jpg
% Wording, abbreviations and medical statements follow the handwritten source.
\ifdefined\settextfont
  \def\BrainTumorsStandalone{0}
\else
  \def\BrainTumorsStandalone{1}
  \documentclass[11pt,a4paper]{article}
  \usepackage[margin=2cm]{geometry}
  \usepackage{amsmath,amssymb,enumitem}
  \usepackage{xepersian}
  \IfFontExistsTF{Amiri}{\settextfont{Amiri}}{\settextfont{FreeSerif}}
  \setlatintextfont{FreeSerif}
  \setlist[itemize]{itemsep=2pt,topsep=3pt,parsep=0pt}
\fi
\providecommand{\en}[1]{\lr{#1}}
\newcommand{\btarrow}{\ensuremath{\leftarrow}}
\newcommand{\btpage}[1]{%
  \clearpage\section*{تومورهای مغزی — تصویر #1}
  {\footnotesize
  \noindent\en{Subject:} ۶ ـ تومورهای مغزی\par
  \noindent\en{Date: \quad / \quad /}\par
  \noindent\en{Sat.\quad Sun.\quad Mon.\quad Tues.\quad Wed.\quad Thu.\quad Fri.}\par}
  \medskip}
\ifnum\BrainTumorsStandalone=1 \begin{document}\fi

\btpage{۱}
\subsection*{مننژیوم‌ها}
\begin{itemize}
  \item اکثراً خوش‌خیم / معمولاً از تحت‌شامه یا آراکنوئید منشأ می‌گیرند.
  \item معمولاً در سینوس ساژیتال، \en{CP angle} (زاویهٔ بین مخچه و پل مغزی)، روی انحناهای نیمکره‌های مغز.
  \item معمولاً در زنان میان‌سال.
  \item در تصویربرداری: یک تودهٔ متصل به دورا که به طور یکنواخت مادهٔ حاجب را جذب می‌کند.
  \item درمان: جراحی و برداشت کامل توده.
\end{itemize}

\subsection*{مدولوبلاستوما}
\begin{itemize}
  \item شایع‌ترین تومور بدخیم مغزی در کودکان.
  \item محل: بخش خلفی ورمیس مخچه / سقف بطن چهارم.
  \item از طریق \en{CSF} منتشر می‌شود.
  \item درمان: جراحی.
  \item پیش‌آگهی قبل از متاستاز خوب است.
\end{itemize}

\subsection*{کرانیوفارنژیوم}
\begin{itemize}
  \item محل: به صورت توده‌ای در ناحیهٔ سوپراسلار.
  \item اغلب کیستیک است.
  \item منشأ \btarrow{} پاچ راتکه / دایورتیکولوم اوروفارنژیال.
  \item \en{MRI} با تزریق: تودهٔ کیستیک سوپراسلار با جذب گادولینیوم.
  \item علائم:
  \begin{itemize}
    \item کودکان: اختلال رشد.
    \item بالغین: اختلال اندوکراین.
    \item نقص بینایی در هر دو.
  \end{itemize}
  \item درمان: جراحی.
\end{itemize}

\btpage{۲}
\subsection*{آستروسایتوما}
\begin{itemize}
  \item شایع‌ترین تومور اولیهٔ مغز.
  \item \en{Grade I}: آستروسایتوم پیلوسیتیک جوانان \btarrow{} پیش‌آگهی پس از جراحی عالی است.
  \item \en{Grade IV}: گلیوبلاستوم مولتی‌فرم \btarrow{} بسیار مهاجم.
  \item \en{Grade I} \btarrow{} شایع در کودکان / اکثراً در مخچه / تومور کیستیک با حاشیهٔ مشخص.
  \item آستروسایتوم \en{High grade} \btarrow{} آستروسایتوم آناپلاستیک و گلیوبلاستوم مولتی‌فرم.
  \begin{itemize}
    \item بیشتر در بزرگ‌سالان / برداشتن کامل تومور امکان‌پذیر نیست.
    \item متاستاز خارج مغزی نادر است.
  \end{itemize}
\end{itemize}

\subsection*{اولیگودندروگلیوما}
\begin{itemize}
  \item اغلب به صورت سوپراتنتوریال / ۳۰\% موارد کلسیفیکاسیون دارند.
  \item \en{MRI} بدون تزریق در \en{$T_2$}: توده‌ای هتروژن / با ادم لوکال و ساب‌کورتیکال ـ همراه کلسیفیکاسیون و تغییرات کیستیک.
  \item از آستروسایتوم خوش‌خیم‌تر است.
  \item درمان: جراحی و برداشتن کامل توده.
\end{itemize}

\subsection*{اپاندیموما}
\begin{itemize}
  \item در کودکان از بطن‌های مغزی / خصوصاً ۴ منشأ می‌گیرد.
  \item تومور با حاشیهٔ مشخص / متاستاز از طریق \en{CSF}.
  \item در صورت برداشتن کامل پیش‌آگهی عالی است.
\end{itemize}

\subsection*{تودهٔ ژرم‌سل}
\begin{itemize}
  \item ژرمینوما شایع‌ترین است که اغلب در دههٔ دوم زندگی رخ می‌دهد.
  \item اکثراً در بطن سوم و نواحی میدلاین خصوصاً ناحیهٔ پینه‌آل.
  \item علائم: اختلال محور هیپوتالاموس ـ هیپوفیز مثل \en{DE}، اختلال میدان بینایی، اختلال خلق، هیدروسفالی.
\end{itemize}

\btpage{۳}
% Continuation of the germ-cell tumor notes from pic02.
\begin{itemize}
  \item در تصویربرداری حاشیهٔ مشخص ندارد / جذب یکنواخت مادهٔ حاجب.
  \item درمان: جراحی.
\end{itemize}

\subsection*{آدنوم هیپوفیز}
\begin{itemize}
  \item تودهٔ خوش‌خیم هیپوفیز قدامی است.
  \item انواع \btarrow{} میکروآدنوم (\en{$<1\,\mathrm{cm}$}) / ماکروآدنوم (\en{$>1\,\mathrm{cm}$}).
  \item علائم: اختلال بینایی، سردرد، دوبینی، رینورهٔ \en{CSF}؛ افزایش یا کاهش هورمون‌های هیپوفیزی.
  \begin{itemize}
    \item کوشینگ / آکرومگالی / آمنوره.
  \end{itemize}
  \item \en{MRI} با و بدون تزریق بهترین روش بررسی توده‌های سلار.
  \item محورهای هورمونی باید بررسی شوند.
  \item درمان: جراحی، رادیوتراپی، هورمون‌درمانی.
\end{itemize}

\subsection*{توده‌های متاستاتیک مغزی}
\begin{itemize}
  \item اکثر موارد از مسیر هماتوژن به مغز می‌رسند.
  \item محل اصلی: ناحیهٔ اتصال مادهٔ سفید به مادهٔ خاکستری.
  \item منشأ شایع: ریه، سپس پستان، \en{GI}، ملانوم.
  \item متاستاز از پروستات، تخمدان و لنفوم هوچکین ـ مغز ناشایع است.
  \item تصویربرداری: توده‌های متعدد با حاشیهٔ نسبتاً مشخص که مادهٔ حاجب را جذب می‌کنند.
\end{itemize}

\subsection*{متاستاز به لپتومننژ (مننژیت کارسینوماتوز)}
\begin{itemize}
  \item آدنوکارسینوم پستان، ریه، \en{GI}، ملانوم.
\end{itemize}

\btpage{۴}
\subsection*{کیست کلوئید}
\begin{itemize}
  \item یک تودهٔ خوش‌خیم با منشأ نامعلوم.
  \item بطن سوم.
  \item علائم: انسداد در مسیر \en{CSF} \btarrow{} هیدروسفالی.
  \item درمان: جراحی.
\end{itemize}

\subsection*{شوانوم}
\begin{itemize}
  \item شایع‌ترین تودهٔ \en{CPA} می‌باشد / خوش‌خیم است.
  \item اکثراً بخش وستیبولار عصب ۸.
  \item درگیری عصب ۵ / درگیری عصب \en{V}.
  \item بیماری همراه: \en{NF-II} \btarrow{} شوانوم شاخهٔ وستیبولار.
  \item علائم بالینی: تودهٔ \en{CPA} به ندرت موجب سرگیجهٔ تنها می‌شود، بلکه با کاهش شنوایی پیش‌روندهٔ یک‌طرفه و وزوز گوش.
  \item شوانوم وستیبولار: بیشتر علائم تعادلی، کاهش شنوایی پیش‌رونده و یک‌طرفه.
  \item[$\ast$] هر بیماری با کاهش شنوایی یک‌طرفهٔ غیرقابل توجیه \btarrow{} \en{MRI} و اودیومتری.
  \item تصویربرداری: توده‌ای که به طور یکنواخت گادولینیوم را جذب می‌کند.
  \item درمان: جراحی.
\end{itemize}

\subsection*{لنفوم اولیهٔ \en{CNS}}
\begin{itemize}
  \item بدخیمی \en{high grade} لنفوسیت‌های \en{B} در \en{CNS}.
  \item شواهدی از وجود لنفوم سیستمیک وجود ندارد.
  \item \en{RF}: نقص ایمنی، دریافت پیوند، \en{AIDS}.
  \item تصویربرداری: در بیمار بدون نقص ایمنی، به صورت توده‌ای با جذب کاملاً یکنواخت مادهٔ حاجب / در بیمار با نقص ایمنی، توده‌های متعدد و با جذب در حاشیهٔ توده (\en{Ring}).
\end{itemize}

\btpage{۵}
% Continuation of primary CNS lymphoma from pic04.
\begin{itemize}
  \item درمان: شیمی‌درمانی.
  \item \en{CS}: پاسخ خوب می‌دهد ولی ظرف چند هفته عود می‌کند.
  \item پیش‌آگهی: خوب نیست.
\end{itemize}

% Handwritten bottom page labels, retained verbatim rather than renumbered:
% pic01: 1; pic02: 3 (also 3 in the lower-left punched hole);
% pic03: 4; pic04: 4 (also 5 in the lower-left punched hole); pic05: 5.
% Printed Date/Subject/day-of-week labels are reproduced in each page header.
\ifnum\BrainTumorsStandalone=1 \expandafter\enddocument\fi
